\documentclass[12pt]{article}
\usepackage[utf8]{inputenc}
\usepackage[T1]{fontenc}
\usepackage{hyperref}
\title{Kombai Language Resources for K-12 Curriculum Development}
\author{Compiled by Claude Sonnet 5 (claude-sonnet-5) for Andr\'as Kornai}
\date{}
\begin{document}
\maketitle

\section{Introduction}
Kombai (ISO 639-3 tyn) is a Papuan language of the Awyu group, spoken in Boven Digoel Regency, South Papua, Indonesia. Wikipedia cites about 5,000 speakers. Joshua Project lists the dialects Central Kombai and Tayan. Kombai should not be confused with Komba (kpf) of Papua New Guinea, which appeared in several search results. Only one substantial scholarly resource turned up.

\section{Primary Educational Resources}
None found.

\section{Linguistic and Reference Materials}
\begin{itemize}
\item Lourens de Vries (1993), \emph{Forms and Functions in Kombai, an Awyu Language of Irian Jaya}, Pacific Linguistics B-108, a grammar (Glottolog bibliographic record): \url{https://glottolog.org/resource/reference/id/96492}
\item Wikipedia article on Kombai: \url{https://en.wikipedia.org/wiki/Kombai_language}
\end{itemize}

\section{Religious and Community Resources}
Joshua Project page on the Kombai in Indonesia, which notes that church services are normally in Indonesian: \url{https://joshuaproject.net/index.php/people_groups/12759}. No Kombai Bible translation was found.

\section{Digital and Audio Resources}
None found. A numerals database entry exists, with a speaker estimate of about 4,000: \url{https://lingweb.eva.mpg.de/channumerals/Kombai.htm}

\section{Curriculum Development Context}
Kombai is not known to be taught in schools, and no dictionary, primer or Scripture was found. The de Vries grammar is the principal basis for any orthography and teaching-grammar work, which would require community and local education authority involvement.
\end{document}